% Gabarit pandoc des corrigés publiés (brut et guidé). Contenu sans marque ni nom de site.
% Compilation : lualatex (latexmk -lualatex).
% Le Markdown reste la source unique : ne pas éditer ce fichier .tex à la main.
\documentclass[11pt,a4paper]{article}

% --- Moteur, polices (Latin Modern) et langue ---------------------------------
\usepackage{iftex}
\ifLuaTeX\else\errmessage{Ce document se compile avec lualatex}\fi
\usepackage{amsmath,amssymb,mathtools}
\usepackage{fontspec}
% mathrm=sym : \mathrm prend ses glyphes dans la police mathématique (ℓ, μ, Ω), comme KaTeX.
\usepackage[math-style=TeX,bold-style=TeX,mathrm=sym]{unicode-math}
% Repli pour les glyphes absents de Latin Modern (indices et exposants Unicode,
% flèches, filets des schémas en texte, symboles) : DejaVu, puis Latin Modern Math
% (ℓ, lettres mathématiques Unicode des légendes).
\directlua{
  luaotfload.add_fallback("repliTexte", {"DejaVuSerif:mode=harf;", "DejaVuSans:mode=harf;", "file:latinmodern-math.otf:mode=harf;"})
  luaotfload.add_fallback("repliMono", {"DejaVuSansMono:mode=harf;", "DejaVuSans:mode=harf;", "file:latinmodern-math.otf:mode=harf;"})
}
\setmainfont{Latin Modern Roman}[RawFeature={fallback=repliTexte}]
\setsansfont{Latin Modern Sans}[RawFeature={fallback=repliTexte}]
\setmonofont{Latin Modern Mono}[RawFeature={fallback=repliMono}]
\setmathfont{Latin Modern Math}
\usepackage[french]{babel}
\frenchsetup{SmallCapsFigTabCaptions=false}

% --- Mise en page --------------------------------------------------------------
\usepackage[a4paper,margin=2.2cm,headheight=14pt,footskip=1.2cm]{geometry}
\usepackage{microtype}
\setlength{\emergencystretch}{4em}
\tolerance=1500
\setlength{\parindent}{0pt}
\setlength{\parskip}{0.55\baselineskip plus 2pt minus 1pt}
\usepackage{xcolor}
% Gris nommés : babel-french rend « ! » actif dans le document.
\definecolor{gris2}{gray}{0.98}
\definecolor{gris3}{gray}{0.97}
\definecolor{gris5}{gray}{0.95}
\definecolor{gris35}{gray}{0.65}
\definecolor{gris40}{gray}{0.60}
\definecolor{gris45}{gray}{0.55}
\definecolor{gris50}{gray}{0.50}
\definecolor{gris55}{gray}{0.45}
\definecolor{gris60}{gray}{0.40}
\definecolor{gris65}{gray}{0.35}
\definecolor{gris70}{gray}{0.30}
\definecolor{gris80}{gray}{0.20}
\definecolor{gris85}{gray}{0.15}
\usepackage{graphicx}
\usepackage{float}
\usepackage{caption}
\captionsetup{font=small,labelformat=empty,justification=centering}
\usepackage{booktabs,longtable,array,multirow,calc}
\usepackage{siunitx}
\sisetup{locale=FR}
\usepackage[version=4]{mhchem}
\usepackage{listings}
\usepackage{fancyhdr}
\usepackage{lastpage}
\usepackage[most]{tcolorbox}
\usepackage{enumitem}
\setlist{nosep,leftmargin=*,topsep=2pt}
\usepackage[hidelinks,unicode]{hyperref}
\usepackage{bookmark}
\hypersetup{pdftitle={Banque PT PT — Français B 2026},pdfauthor={Rédigé par IA},pdfcreator={LaTeX},pdfsubject={Corrigé guidé}}
\usepackage{xurl}

% --- Pandoc : listes, tableaux, figures ----------------------------------------
\providecommand{\passthrough}[1]{#1}
\providecommand{\tightlist}{\setlength{\itemsep}{0pt}\setlength{\parskip}{0pt}}

\makeatletter
\def\maxwidth{\ifdim\Gin@nat@width>\linewidth\linewidth\else\Gin@nat@width\fi}
\def\maxheight{\ifdim\Gin@nat@height>0.55\textheight 0.55\textheight\else\Gin@nat@height\fi}
\makeatother
\setkeys{Gin}{width=\maxwidth,height=\maxheight,keepaspectratio}
\setcounter{secnumdepth}{-\maxdimen}
\providecommand{\pandocbounded}[1]{#1}
\usepackage{longtable,booktabs,array}

% --- Code (listings) -----------------------------------------------------------
\lstset{
  basicstyle=\ttfamily\small,
  columns=fullflexible,
  keepspaces=true,
  breaklines=true,
  breakatwhitespace=false,
  postbreak=\mbox{\textcolor{gris50}{$\hookrightarrow$}\space},
  frame=single,
  rulecolor=\color{gris35},
  framesep=4pt,
  xleftmargin=5pt,xrightmargin=5pt,
  backgroundcolor=\color{gris3},
  keywordstyle=\bfseries,
  commentstyle=\color{gris65},
  stringstyle=\color{gris80},
  showstringspaces=false,
  tabsize=4,
  extendedchars=true,
  aboveskip=6pt,belowskip=6pt,
}

% --- Macros des corrigés (KaTeX ↔ LaTeX) ---------------------------------------
% Les corrigés sont écrits pour KaTeX : on définit ici ce que KaTeX connaît et
% que LaTeX n'a pas, sans jamais retoucher le Markdown.
\providecommand{\R}{\mathbb{R}}
\providecommand{\N}{\mathbb{N}}
\providecommand{\Z}{\mathbb{Z}}
\providecommand{\Q}{\mathbb{Q}}
\providecommand{\C}{\mathbb{C}}
\providecommand{\dd}{\mathrm{d}}
\providecommand{\llbracket}{\lBrack}
\providecommand{\rrbracket}{\rBrack}
\providecommand{\lBrace}{\lbrace}
\providecommand{\rBrace}{\rbrace}
\providecommand{\Box}{\square}
\providecommand{\Diamond}{\lozenge}
\providecommand{\varnothing}{\emptyset}
\providecommand{\argmax}{\operatorname*{arg\,max}}
\providecommand{\argmin}{\operatorname*{arg\,min}}
\renewcommand{\leq}{\leqslant}
\renewcommand{\geq}{\geqslant}

% Équation hors texte : réduite à la largeur de la ligne si elle déborde.
\newsavebox{\boiteequation}
\newcommand{\equationajustee}[1]{%
  \sbox{\boiteequation}{$\displaystyle #1$}%
  \ifdim\wd\boiteequation>\linewidth
    \[\resizebox{\linewidth}{!}{\usebox{\boiteequation}}\]%
  \else
    \[#1\]%
  \fi}

% --- Encadrés pédagogiques (sobres, lisibles en noir et blanc) -----------------
\tcbset{
  encadre/.style={
    enhanced,breakable,sharp corners,
    colback=white,colframe=gris55,
    boxrule=0.4pt,leftrule=2.2pt,
    left=6pt,right=6pt,top=3pt,bottom=3pt,
    before skip=6pt,after skip=6pt,
    fonttitle=\bfseries\small,coltitle=black,
    attach title to upper,after title={\quad},
    parbox=false,
  }
}
\newtcolorbox{indication}[1]{encadre,colframe=gris45,colback=gris2,title={Indication~#1}}
\newtcolorbox{methode}{encadre,colframe=gris70,title={Méthode}}
\newtcolorbox{piege}{encadre,colframe=gris85,boxrule=0.8pt,leftrule=2.2pt,title={Piège}}
\newtcolorbox{jury}{encadre,colframe=gris60,colback=gris5,title={Jury}}
\newtcolorbox{encadrenote}{encadre,colframe=gris40}

% --- Titres de section ---------------------------------------------------------
\usepackage{titlesec}
\titleformat{\section}{\Large\bfseries}{}{0pt}{}
\titleformat{\subsection}{\large\bfseries}{}{0pt}{}
\titleformat{\subsubsection}{\normalsize\bfseries}{}{0pt}{}
\titlespacing*{\section}{0pt}{14pt plus 4pt minus 2pt}{6pt}
\titlespacing*{\subsection}{0pt}{10pt plus 3pt minus 2pt}{4pt}
\titlespacing*{\subsubsection}{0pt}{8pt plus 2pt minus 2pt}{2pt}

% --- En-tête et pied de page ---------------------------------------------------
\pagestyle{fancy}
\fancyhf{}
\fancyhead[L]{\small\textsc{Corrigé guidé} · Banque PT PT}
\fancyhead[R]{\small Français B 2026}
\fancyfoot[L]{\footnotesize Corrigé rédigé par IA, non officiel}
\fancyfoot[R]{\footnotesize \thepage\,/\,\pageref*{LastPage}}
\renewcommand{\headrulewidth}{0.4pt}
\renewcommand{\footrulewidth}{0pt}
\fancypagestyle{plain}{\fancyhf{}\fancyfoot[R]{\footnotesize \thepage\,/\,\pageref*{LastPage}}\fancyfoot[L]{\footnotesize Corrigé rédigé par IA, non officiel}\renewcommand{\headrulewidth}{0pt}}

\begin{document}
\thispagestyle{plain}
\begin{center}
  {\small\scshape Corrigé guidé}\par\smallskip
  {\LARGE\bfseries Banque PT PT — Français B 2026\par}\medskip
  {\large Banque PT · PT · Français B · 2026\par}
\end{center}
\vspace{2pt}\hrule\vspace{8pt}

\emph{Corrigé rédigé par IA, non officiel. Signalez toute erreur.}

\hypertarget{pruxe9sentation-du-sujet}{%
\section{Présentation du sujet}\label{pruxe9sentation-du-sujet}}

\textbf{L'épreuve.} Quatre heures, sans document ni dictionnaire ni appareil électronique, pour deux exercices notés séparément : un \textbf{résumé} de texte (8 points) en 200 mots, avec une marge de ± 10 \% (180 à 220 mots), une barre oblique tous les vingt mots et le total inscrit à la fin ; une \textbf{dissertation} (12 points) sur une citation tirée du même texte, à traiter à partir des trois œuvres au programme. Les « questions » de ce corrigé guidé sont les étapes du travail : trois pour le résumé, cinq pour la dissertation.

\textbf{Le texte.} L'extrait de François Dagognet (\emph{Considérations sur l'idée de nature}, Vrin, 2000) compte environ 1 400 mots : le résumé opère donc une réduction au septième. C'est un texte polémique, à thèse frontale, écrit à la première personne du pluriel. Son mouvement :

\begin{longtable}[]{@{}
  >{\raggedright\arraybackslash}p{(\columnwidth - 4\tabcolsep) * \real{0.1390}}
  >{\raggedright\arraybackslash}p{(\columnwidth - 4\tabcolsep) * \real{0.2677}}
  >{\raggedright\arraybackslash}p{(\columnwidth - 4\tabcolsep) * \real{0.5633}}@{}}
\toprule\noalign{}
\begin{minipage}[b]{\linewidth}\raggedright
Mouvement
\end{minipage} & \begin{minipage}[b]{\linewidth}\raggedright
Bornes
\end{minipage} & \begin{minipage}[b]{\linewidth}\raggedright
Idée directrice
\end{minipage} \\
\midrule\noalign{}
\endhead
\bottomrule\noalign{}
\endlastfoot
1. L'alternative & « Si nul ne peut nier\ldots{} le plus d'attrait. » & Face aux dégâts de l'industrie, deux remèdes : supprimer la cause (écologisme, arrêt de la croissance, énergies douces) ou corriger les effets (éco-industrie, « industrie de l'industrie »). L'auteur choisit la seconde, tout en reconnaissant que la première séduit davantage. \\
2. Le déplacement & « Le problème relève d'ailleurs\ldots{} en alliage léger » (O. Vallet) & Le problème est philosophique : il tient à notre représentation de la nature. Preuve : la publicité vend « frauduleusement » du naturel (fruits, œufs, blé dur) ; le ski de fond, signe du retour à la nature, est mécanique et chimique de part en part. \\
3. L'explication & « Pourquoi cette préférence\ldots{} (le naturalisme). » & On refuse l'imitation, qui déforme ou décalque et déçoit toujours ; d'où la victoire du « fruste ». On oublie que tout est transformé, jusqu'à nos jardins peuplés de plantes exotiques. \\
4. Kant & « Kant lui-même\ldots{} (du fait de la trahison). » & La fleur artificielle et le faux rossignol : la morale s'immisce dans le jugement esthétique et condamne le leurre comme un mensonge. Double faute de l'imitation : elle mécanise la nature et la scinde (le chant sans l'oiseau). Désenchantement ontologique et moral. \\
5. Le diagnostic & « Mais la rhétorique publicitaire\ldots{} des siècles ? » & La publicité use d'un lexique mythique (pureté, simplicité, pérennité) ; l'acheteur obéit moins à sa santé qu'à une « religiosité naïve » de l'incréé, du premier, du sauvage, reçu comme un don du ciel ; il reproche à l'usine de briser les complexes naturels et d'aller trop vite. \\
6. La réfutation & « À l'opposé de ces croyances\ldots{} s'acidifie. » & L'eau de source (chère à Rousseau) est dangereuse et doit être traitée ; l'eau minérale suppose industrie et commerce, et « naturelle » est un jeu de mots ; le sol non travaillé se tasse et s'appauvrit. \\
7. La conclusion & « Partout, il faut intervenir\ldots{} ``âge d'or'' ! » & Il faut intervenir, transformer de fond en comble, et renoncer à l'illusion d'un âge d'or. \\
\end{longtable}

Le point capital est la \textbf{polarité} du texte : ce n'est pas une dénonciation de l'industrie, c'est la dénonciation de l'illusion naturaliste par un partisan de l'industrie corrigée. L'inverser est un contresens global.

\textbf{La dissertation.} Le libellé reprend la phrase centrale du mouvement 5 et demande, au conditionnel, si cette « religiosité naïve » \emph{caractérise} le rapport de l'homme à la nature. Les œuvres au programme s'y prêtent remarquablement :

\begin{itemize}
\tightlist
\item
  \textbf{Canguilhem, \emph{La Connaissance de la vie}} (périmètre : Introduction « La pensée et le vivant » ; « L'expérimentation en biologie animale » ; « Machine et organisme », « Le vivant et son milieu », « Le normal et le pathologique », « La monstruosité et le monstrueux ») : la connaissance comme analyse, la technique comme phénomène vital, le milieu propre du vivant, la santé comme normativité, le monstre et le monstrueux.
\item
  \textbf{Verne, \emph{Vingt mille lieues sous les mers}} : Nemo et son culte de la mer libre, le \emph{Nautilus} comme artefact total nourri de l'océan, la nomenclature d'Aronnax et de Conseil, les dangers marins, les cachalots, l'éperonnage final.
\item
  \textbf{Haushofer, \emph{Le Mur invisible}} : la nature de villégiature devenue prison, la survie sur les réserves du monde disparu, le travail paysan, l'attachement aux bêtes, le meurtre final.
\end{itemize}

\textbf{Difficulté et gestion du temps.} Le résumé est très accessible à condition de respecter strictement le format et l'ordre du texte ; ce sont des points qu'il ne faut pas laisser passer. La dissertation est plus exigeante : le sujet-citation piège à la fois la paraphrase du texte et la réfutation sommaire. Répartition indicative : 1 h 10 pour le résumé (lecture, découpage, brouillon, décompte), 2 h 40 pour la dissertation (30 min d'analyse et de plan, 2 h de rédaction), 10 min de relecture.

\textbf{Ce qu'en dit le jury.} Le rapport 2026 compte 2707 candidats, pour une moyenne de 10,01 et un écart type de 3,99 (p.~1). Le résumé obtient en moyenne 4,55 sur 8 (p.~3), la dissertation 6,15 sur 12 (p.~5). Trois points ont fait la différence :

\begin{itemize}
\tightlist
\item
  le \textbf{décompte des mots}, vérifié sur toutes les copies : 219 copies y ont perdu de 1 à 4 points (p.~3) ;
\item
  la \textbf{seconde partie du texte} (les raisons du culte du naturel), mal articulée dans beaucoup de résumés (p.~4) ;
\item
  la \textbf{pertinence} de la dissertation, que le jury qualifie de « critère très discriminant » (p.~5), et en particulier la prise au sérieux du mot « religiosité » (p.~6).
\end{itemize}

\begin{jury}

Le rapport 2026 déconseille de commencer par la dissertation : « C'est là un mauvais calcul. En effet, la bonne intelligence du sujet sera toujours favorisée par une compréhension fine de l'ensemble dont il est extrait, ce que permet précisément le résumé. » (p.~4)

\end{jury}

\hypertarget{partie-i-le-ruxe9sumuxe9}{%
\section{Partie I --- Le résumé}\label{partie-i-le-ruxe9sumuxe9}}

Trois étapes : comprendre la logique du texte, la hiérarchiser pour répartir les mots, puis rédiger et compter. Outils : repérage des connecteurs, reformulation, décompte rigoureux.

\hypertarget{uxe9tape-1-lire-le-texte-et-duxe9gager-sa-thuxe8se}{%
\subsection{Étape 1 --- Lire le texte et dégager sa thèse}\label{uxe9tape-1-lire-le-texte-et-duxe9gager-sa-thuxe8se}}

\begin{indication}{1}

Cherchez d'abord la phrase où l'auteur prend parti à la première personne : elle donne la polarité de tout le texte.

\end{indication}

\begin{indication}{2}

Les exemples (publicité, ski de fond, Kant, eau de source) ne sont pas juxtaposés : demandez-vous, pour chacun, s'il prouve un constat, explique une cause ou réfute une croyance.

\end{indication}

\hypertarget{solution}{%
\subsubsection{Solution}\label{solution}}

La thèse est énoncée au premier paragraphe : « Nous défendons assurément cette seconde solution », celle de l'éco-industrie qui corrige les effets de l'industrialisation, contre celle qui en supprime la cause. Tout le reste du texte explique pourquoi la première solution séduit malgré tout, et pourquoi cette séduction est une illusion.

L'argumentation suit un ordre précis : un \textbf{déplacement} (le problème est philosophique, il tient à notre représentation de la nature), deux \textbf{preuves} (la publicité du naturel, le ski de fond), une \textbf{explication} (le refus de l'imitation, illustré et confirmé par Kant), un \textbf{diagnostic} (une « religiosité naïve » de l'originel), une \textbf{réfutation} (l'eau, le sol) et une \textbf{conclusion} impérative (« Partout, il faut intervenir »). La problématique du texte peut se formuler ainsi : faut-il, pour remédier aux dégâts de l'industrie, revenir à une nature intacte, ou dissiper l'illusion qui fait préférer le donné au transformé ?

\begin{piege}

Lire le texte comme une satire de la publicité ou, pire, comme un plaidoyer écologiste. La publicité n'est qu'une preuve ; la cible est la croyance au « naturel », et l'auteur défend l'intervention technique.

\end{piege}

\begin{jury}

Selon le rapport 2026, les premières lignes du texte ont été « généralement bien comprises », mais « la seconde partie du texte qui tente d'expliquer les raisons d'une attirance et d'une fascination collective pour le naturel est plus confusément perçue » (p.~4). C'est là que se gagnent les points.

\end{jury}

\begin{piege}

Le contresens sur Kant. Le jury a lu que, selon Kant, « les hommes se sentent dupés lorsqu'ils découvrent des fleurs artificielles qui imitent le chant du rossignol » (p.~4), et même que Kant avait lutté contre le réchauffement climatique. Kant analyse deux leurres distincts, la fleur artificielle et le garçon caché qui siffle ; Dagognet s'en sert pour expliquer la défiance envers l'imitation, non pour parler d'écologie.

\end{piege}

\hypertarget{uxe9tape-2-duxe9couper-le-texte-et-ruxe9partir-les-mots}{%
\subsection{Étape 2 --- Découper le texte et répartir les mots}\label{uxe9tape-2-duxe9couper-le-texte-et-ruxe9partir-les-mots}}

\begin{indication}{1}

Découpez par unités de sens, non par paragraphes : un même paragraphe peut contenir deux étapes, et deux paragraphes peuvent n'en former qu'une.

\end{indication}

\begin{indication}{2}

Répartissez les 200 mots selon le poids argumentatif de chaque unité, non selon sa longueur : un exemple pittoresque mais secondaire se résume en une proposition.

\end{indication}

\hypertarget{solution-1}{%
\subsubsection{Solution}\label{solution-1}}

Le découpage en sept mouvements de la présentation donne la répartition suivante, pour un total visé d'environ 205 mots (au centre de la fourchette) :

\begin{longtable}[]{@{}
  >{\raggedright\arraybackslash}p{(\columnwidth - 4\tabcolsep) * \real{0.2819}}
  >{\raggedright\arraybackslash}p{(\columnwidth - 4\tabcolsep) * \real{0.1221}}
  >{\raggedright\arraybackslash}p{(\columnwidth - 4\tabcolsep) * \real{0.5660}}@{}}
\toprule\noalign{}
\begin{minipage}[b]{\linewidth}\raggedright
Mouvement
\end{minipage} & \begin{minipage}[b]{\linewidth}\raggedright
Budget
\end{minipage} & \begin{minipage}[b]{\linewidth}\raggedright
Ce qu'il faut garder
\end{minipage} \\
\midrule\noalign{}
\endhead
\bottomrule\noalign{}
\endlastfoot
1. L'alternative et le choix & ≈ 40 & Les deux remèdes, le choix de l'auteur, l'attrait de l'autre voie \\
2. Déplacement et preuves & ≈ 40 & Le problème est philosophique ; publicité et ski de fond réunis en une phrase \\
3. Première raison : l'imitation & ≈ 15 & Refus de l'imitation ; oubli que tout est déjà transformé \\
4. Kant & ≈ 25 & Le leurre condamné comme mensonge ; mécaniser et mutiler ; déception ontologique et morale \\
5. Deuxième et troisième raisons & ≈ 45 & Lexique mythique et dévotion à l'originel reçu du ciel ; reproche fait à l'usine, trop rapide pour égaler la lenteur de la nature \\
6--7. Réfutation et conclusion & ≈ 40 & Eau de source, eau minérale, sol ; il faut intervenir ; pas d'âge d'or \\
\end{longtable}

Les détails (œufs, blé dur, Vallet, la gourde, l'iode, \emph{silvaticus}) disparaissent : ce sont des illustrations, pas des articulations.

\begin{methode}

Le budget de mots se fixe au brouillon, avant toute rédaction. On écrit ensuite chaque unité un peu trop longue, puis on élague : c'est plus rapide que de chercher le compte juste du premier coup.

\end{methode}

\begin{jury}

Le plan du texte proposé par le rapport 2026 (p.~3) a trois temps : le problème posé par l'industrialisation (§ 1-2), les raisons de la promotion du naturel (§ 3-5), la nécessité de l'action transformatrice (§ 6). Les raisons sont au nombre de trois : la méfiance envers l'imitation, à laquelle se rattache Kant ; la valorisation mythique de l'âge d'or ; enfin « la mise en cause des rythmes industriels qui ne respectent pas le temps long de la nature ». Un résumé qui oublie cette troisième raison perd une articulation.

\end{jury}

\begin{piege}

Sacrifier le premier mouvement, jugé abstrait, pour commencer par la publicité : le résumé perd alors la thèse et devient une simple satire du marketing.

\end{piege}

\hypertarget{uxe9tape-3-ruxe9diger-compter-pruxe9senter}{%
\subsection{Étape 3 --- Rédiger, compter, présenter}\label{uxe9tape-3-ruxe9diger-compter-pruxe9senter}}

\begin{indication}{1}

Gardez l'ordre du texte et son énonciation : « nous » pour la thèse de l'auteur, « on » pour la position qu'il critique ; aucun commentaire du type « l'auteur montre que ». Les trois raisons du culte du naturel doivent être reliées par des connecteurs qui les comptent.

\end{indication}

\begin{indication}{2}

Reformulez au lieu de recopier, en particulier la formule « l'incréé, le premier et le sauvage », qui est réservée à la dissertation. Fixez votre règle de décompte (apostrophe, trait d'union) et appliquez-la jusqu'au bout.

\end{indication}

\hypertarget{solution-2}{%
\subsubsection{Solution}\label{solution-2}}

Convention de décompte : est compté pour un mot tout ensemble de lettres séparé du suivant par une espace ou une apostrophe ; un mot composé à trait d'union compte pour un mot. Une barre oblique marque chaque vingtième mot ; le total cumulé suit entre parenthèses.

\begin{quote}

Nul ne conteste les ravages d'une industrialisation débridée, mais deux remèdes s'opposent : abolir leur cause, par la décroissance / (20) ou les énergies douces, ou corriger leurs effets par une industrie réparatrice. Nous retenons la seconde voie, quoique la première / (40) séduise davantage.

La question est philosophique : tout dépend de notre idée de la nature. La publicité l'atteste, qui vante / (60) comme naturels des aliments issus d'engrais et d'élevages forcés, ou le ski de fond, tout mécanique et chimique. / (80)

Pourquoi cet engouement ? D'abord, on rejette l'imitation, toujours décevante, en oubliant que tout est déjà transformé. Kant lui-même, / (100) devant un faux rossignol, dénonçait un mensonge : imiter mécanise la nature et la mutile, d'où une déception ontologique et / (120) morale.

Ensuite, flatté par un lexique mythique, l'acheteur obéit moins au souci de sa santé qu'à une dévotion / (140) crédule envers l'originel, reçu du ciel. Enfin, il soupçonne l'usine de briser l'unité du donné et de / (160) bâcler ce que la terre mûrit lentement.

Or la nature n'est ni pure ni sûre : l'eau de source / (180) doit être assainie, l'eau minérale exige machines et commerce, la terre laissée à elle-même s'épuise. Il faut partout / (200) intervenir : renonçons au mythe d'un âge d'or.

\textbf{209 mots}

\end{quote}

Contrôle : 209 mots avec la convention annoncée ; 211 si le trait d'union sépare aussi ; 190 si l'apostrophe ne sépare pas. Le résumé est donc dans la fourchette (180 à 220) quelle que soit la règle appliquée par le correcteur.

La fourchette qui fait foi est celle du sujet, de 180 à 220 mots. Le rapport 2026 écrit « entre 180 et 200 » (p.~3), mais sa propre proposition de résumé compte 218 mots : il s'agit d'une coquille.

\begin{methode}

Viser le centre de la fourchette, et non ses bords : à 218 mots, un désaccord de convention suffit à faire sortir du format. Vérifier le décompte en recomptant chaque tranche de vingt mots, barre par barre.

\end{methode}

\begin{piege}

Annoncer un total faux. Selon le rapport 2026, « Toutes les copies font l'objet d'une vérification systématique sur ce point » ; 219 copies, qui falsifiaient leur décompte, ont perdu de 1 à 4 points (p.~3).

\end{piege}

\begin{jury}

Le jury « a par ailleurs apprécié les copies qui ont précisé pour chaque tranche de 20 mots le total cumulé (20, puis 40, 60 etc.) » (p.~3). Il attend un résumé « fidèle au circuit argumentatif logique observé par le texte », qui rende la progression perceptible « au moyen de paragraphes différenciés et de connecteurs logiques » (p.~4).

\end{jury}

\begin{piege}

Chercher à tout prix des synonymes de « nature » ou d'« industrialisation ». Le rapport 2026 rappelle que la reformulation porte sur les idées et que ces termes, « au cœur du texte, pouvaient difficilement être remplacés » (p.~4).

\end{piege}

\hypertarget{partie-ii-la-dissertation}{%
\section{Partie II --- La dissertation}\label{partie-ii-la-dissertation}}

Cinq étapes : analyser le libellé, en tirer une problématique, construire un plan qui distribue les trois œuvres, rédiger l'introduction, puis le développement et la conclusion. L'enjeu est de discuter la thèse de Dagognet sans la paraphraser ni la balayer.

\hypertarget{uxe9tape-4-analyser-le-libelluxe9}{%
\subsection{Étape 4 --- Analyser le libellé}\label{uxe9tape-4-analyser-le-libelluxe9}}

\begin{indication}{1}

Distinguez ce qui est posé (la citation) de ce qui est demandé (le verbe au conditionnel). Que veut dire \emph{caractériser} un rapport ?

\end{indication}

\begin{indication}{2}

Pesez chaque terme de la citation : le suffixe de \emph{religiosité}, le sens de \emph{naïve}, les trois termes de la série finale, et l'adjectif qui qualifie les « combinaisons humaines ».

\end{indication}

\hypertarget{solution-3}{%
\subsubsection{Solution}\label{solution-3}}

\begin{longtable}[]{@{}
  >{\raggedright\arraybackslash}p{(\columnwidth - 2\tabcolsep) * \real{0.3135}}
  >{\raggedright\arraybackslash}p{(\columnwidth - 2\tabcolsep) * \real{0.6565}}@{}}
\toprule\noalign{}
\begin{minipage}[b]{\linewidth}\raggedright
Terme
\end{minipage} & \begin{minipage}[b]{\linewidth}\raggedright
Ce qu'il engage
\end{minipage} \\
\midrule\noalign{}
\endhead
\bottomrule\noalign{}
\endlastfoot
« le rapport de l'homme à la nature » & Non la nature en soi, mais une relation : ce que l'homme fait de la nature dit aussi ce qu'il croit être. \\
« se caractériserait-il » & \emph{Caractériser}, c'est définir par un trait propre et essentiel, non par un trait parmi d'autres. Le conditionnel met la thèse à distance : il faut la juger, non l'illustrer. \\
« une religiosité naïve » & \emph{Religiosité} et non religion : un sentiment diffus du sacré, sans dogme ni examen ; le terme est dépréciatif. \emph{Naïve} : crédule, ignorante de ses propres conditions. \\
« préfère » & Un acte de valorisation, non de connaissance : le rapport visé relève du goût et de la croyance. \\
« aux combinaisons humaines douteuses » & Le composé, le mélangé, le fabriqué, imputable à quelqu'un, et suspect à la fois de fausseté et de faute. \\
« l'incréé, le premier et le sauvage » & Trois registres : théologique (ce qui n'a pas d'auteur), chronologique (ce qui précède toute intervention), culturel (\emph{silva}, la forêt : le non-domestiqué). \\
\end{longtable}

Le libellé suppose donc deux choses : que le partage entre le naturel et l'artificiel ait un sens, et que la préférence pour le premier soit une illusion. C'est ce double présupposé qu'il faudra éprouver.

\begin{piege}

Disserter sur « l'homme et la nature » en général, ou sur l'écologie, en oubliant que le sujet porte sur la validité d'une caractérisation précise. Le rapport 2026 vise surtout le plan qui oppose les bons et les mauvais aspects de la nature : « Car c'est bien le rapport de l'homme à la nature qu'il s'agissait ici d'interroger, et non, simplement les bienfaits ou méfaits de celle-ci. » (p.~6)

\end{piege}

\begin{jury}

Le jury relève la tendance à aplatir le sujet, « notamment en ne tenant aucun compte du sens du mot ``religiosité'' qui engageait un rapport à la nature qui serait chargé de sacré et de mysticisme » (p.~6). Il prévient aussi que l'analyse « ne saurait consister en une liste de synonymes » : « Le jury attend avant tout une appréhension globale du sens du sujet, qui servira de fondement à la discussion. » (p.~6) Le tableau ci-dessus est un outil de brouillon ; dans la copie, l'analyse se rédige en un paragraphe qui conduit au problème.

\end{jury}

\hypertarget{uxe9tape-5-probluxe9matiser}{%
\subsection{Étape 5 --- Problématiser}\label{uxe9tape-5-probluxe9matiser}}

\begin{indication}{1}

Essayez les deux réponses immédiates, « oui » et « non », et cherchez pourquoi chacune est intenable.

\end{indication}

\begin{indication}{2}

Demandez-vous si la préférence pour l'originel ne pourrait pas être le signe d'autre chose que d'une ignorance : que cherche-t-on vraiment quand on préfère le chant du rossignol à celui de son imitateur ?

\end{indication}

\hypertarget{solution-4}{%
\subsubsection{Solution}\label{solution-4}}

Répondre « oui » revient à paraphraser Dagognet et à réduire tout attachement à la nature à une superstition, ce que dément l'expérience. Répondre « non » sans précaution revient à défendre l'âge d'or que le texte réfute avec de bons arguments (l'eau qui contamine, le sol qui s'épuise). La difficulté est donc réelle : la préférence pour l'originel est à la fois une illusion et quelque chose de plus qu'une illusion.

\textbf{Problématique} : la préférence pour l'incréé, le premier et le sauvage est-elle le trait qui définit notre rapport à la nature, ou le symptôme confus d'une exigence de valeur qu'un rapport purement opératoire ne sait pas formuler ?

\begin{methode}

Une problématique n'est pas la reprise interrogative du libellé : elle énonce la difficulté qui oblige à dépasser les réponses immédiates, et elle laisse déjà deviner le plan.

\end{methode}

\begin{jury}

Pour le rapport 2026, la problématique pouvait interroger le bien-fondé de l'affirmation, « mais aussi, au-delà, en interrogeant le bienfondé de cette dissociation discutable, en particulier entre la nature et le sujet humain pensant » (p.~5). C'est ce que fait la problématique retenue, en se demandant si l'opposition de l'originel et du fabriqué tient.

\end{jury}

\begin{piege}

Mêler interrogation directe et indirecte. Le jury cite « on se demandera si les hommes préfèrent-ils le sauvage ? » (p.~2) : il faut écrire « on se demandera si les hommes préfèrent le sauvage », sans inversion ni point d'interrogation.

\end{piege}

\hypertarget{uxe9tape-6-construire-le-plan-et-distribuer-les-ux153uvres}{%
\subsection{Étape 6 --- Construire le plan et distribuer les œuvres}\label{uxe9tape-6-construire-le-plan-et-distribuer-les-ux153uvres}}

\begin{indication}{1}

Un plan dialectique convient : concéder, contester, déplacer. Mais la troisième partie doit changer la question, non proposer un « équilibre » entre nature et technique.

\end{indication}

\begin{indication}{2}

Cherchez dans Canguilhem le concept qui permet de juger une intervention sans invoquer son origine : il se trouve dans l'un des chapitres de la partie « Philosophie ».

\end{indication}

\begin{indication}{3}

Vérifiez au brouillon, par un tableau œuvres × parties, que chaque œuvre apparaît dans chaque partie et qu'elle y fait avancer l'argument.

\end{indication}

\hypertarget{solution-5}{%
\subsubsection{Solution}\label{solution-5}}

\textbf{I. Une religiosité de l'originel travaille bien notre rapport à la nature, et elle est naïve} (concession)
- A. Le prestige de l'intact : Nemo et le culte de la mer libre ; la montagne de villégiature du \emph{Mur invisible} ; l'horreur sacrée du monstre chez Canguilhem.
- B. L'oubli des médiations : le \emph{Nautilus}, artefact total ; la survie sur les réserves du monde disparu ; l'histoire de la notion de milieu.
- C. Une nature sans bienveillance : la vallée indifférente, la mer qui emprisonne et tue, la vie capable de manquer ses formes ; la crainte, autre face du sacré.

\textbf{II. Ce n'est pourtant pas le trait qui définit notre rapport à la nature : celui-ci est d'abord opératoire} (contestation)
- A. Connaître : Aronnax et Conseil ; la connaissance comme analyse et ses limites (Introduction, « L'expérimentation en biologie animale »).
- B. Transformer : « Machine et organisme », la technique comme phénomène vital ; le \emph{Nautilus} nourri de la mer ; le travail paysan de la narratrice.
- C. Habiter : \emph{Umwelt} et \emph{Umgebung} (« Le vivant et son milieu ») ; la vitre du salon ; la vallée devenue monde vécu.

\textbf{III. Le vrai partage n'est pas entre le donné et le fabriqué, mais entre des normes de vie} (déplacement)
- A. Ce que la préférence pour le donné vise : l'autonomie du vivant (le rossignol de Kant), la liberté (Nemo), la responsabilité (les bêtes de la narratrice).
- B. Du naturel à la norme : « Le normal et le pathologique » ; la précarité de la narratrice ; l'arbitraire de Nemo (baleines épargnées, cachalots massacrés).
- C. Le désastre vient de la violence, non de la technique : l'homme à la hache ; l'éperonnage du navire de guerre ; le mur soupçonné d'être une arme ; un savoir enraciné dans la vie.

\begin{methode}

La transition du II au III est fournie par une insuffisance : si tout rapport à la nature est intervention, qu'est-ce qui distingue une intervention qui fait vivre d'une intervention qui détruit ? Une transition argumentée vaut mieux qu'un connecteur décoratif.

\end{methode}

\begin{piege}

Organiser une partie par œuvre (« d'abord Canguilhem, puis Verne, enfin Haushofer ») : on obtient trois résumés successifs et aucun raisonnement. Autre écueil : traiter Canguilhem comme un réservoir de thèses et les romans comme un réservoir d'anecdotes, alors que les romans portent aussi des idées.

\end{piege}

\begin{jury}

Le rapport 2026 propose un plan en trois temps proche de celui-ci (p.~5) : sacralisation de la nature brute ; réhabilitation des artefacts devant une nature hostile ; une culture qui permet de retrouver la nature, « cette fois sans ``religiosité naïve'' ». Il accepte deux parties : « Deux parties qui affrontent véritablement le sujet peuvent donner lieu à une très bonne note. » (p.~6) Pour les exemples, il s'est satisfait « de deux exemples valables pris dans deux œuvres différentes ; la référence à une seule des trois œuvres ne saurait suffire » (p.~6). Il pénalise les copies « qui ont escamoté l'œuvre philosophique » (p.~6).

\end{jury}

\begin{piege}

Oublier la technique. Le jury signale « l'absence de prise en compte de la dualité entre nature et technique », la technique étant « peu envisagée en tant qu'objet de défiance elle aussi, ou encore en tant que moyen de se prémunir contre les menaces de la nature » (p.~6). Le plan retenu la traite sous ses trois faces : moyen de survie (I.B, II.B), forme de vie (II.B), objet de crainte légitime (III.C).

\end{piege}

\hypertarget{uxe9tape-7-ruxe9diger-lintroduction}{%
\subsection{Étape 7 --- Rédiger l'introduction}\label{uxe9tape-7-ruxe9diger-lintroduction}}

\begin{indication}{1}

Quatre temps : une entrée en matière qui pose déjà le problème, la citation et l'analyse de ses termes, la problématique, l'annonce du plan.

\end{indication}

\begin{indication}{2}

Pour l'accroche, préférez une scène tirée d'une œuvre au programme à un exemple repris du texte résumé : vous évitez le soupçon de paraphrase.

\end{indication}

\hypertarget{solution-6}{%
\subsubsection{Solution}\label{solution-6}}

Lorsque la narratrice du \emph{Mur invisible} arrive au pavillon de chasse de ses cousins, elle vient goûter, le temps d'un week-end, une montagne qu'on aime parce qu'elle est intacte. Quelques heures plus tard, une paroi transparente la sépare du reste des hommes, et cette nature de villégiature devient l'unique monde où elle devra survivre. Le roman de Marlen Haushofer commence ainsi là où se termine une certaine idée de la nature : celle d'un sanctuaire que l'on visite et auquel on prête une pureté que rien de fabriqué n'égalerait.

C'est cette idée que François Dagognet, dans ses \emph{Considérations sur l'idée de nature}, attaque de front. Le consommateur qui recherche le « naturel » obéirait, selon lui, à « une religiosité naïve » qui « préfère aux combinaisons humaines douteuses l'incréé, le premier et le sauvage ». La \emph{religiosité} n'est pas une religion : c'est une dévotion diffuse, sans dogme ni examen, qui entoure un objet d'une aura sacrée ; elle est dite \emph{naïve} parce qu'elle ignore ses propres conditions. Son objet se décline en trois termes : l'\emph{incréé}, qui ne doit rien à aucun auteur ; le \emph{premier}, qui précède toute intervention ; le \emph{sauvage}, qui pousse, selon l'étymologie rappelée par l'auteur, hors de la forêt cultivée. Face à lui, les « combinaisons humaines » --- le composé, le mélangé, le fabriqué --- sont tenues pour « douteuses », suspectes à la fois de fausseté et de faute. Le sujet, au conditionnel, demande si ce trait \emph{caractérise} le rapport de l'homme à la nature, c'est-à-dire s'il en est le propre.

La question est difficile parce que les deux réponses immédiates échouent. Souscrire à la thèse, c'est réduire tout attachement à la nature à une superstition, ce que dément l'expérience : on ne s'attache pas à un paysage comme on adore une idole. La rejeter sans plus, c'est réhabiliter l'« âge d'or » que Dagognet réfute avec de bons arguments. D'où notre problème : la préférence pour l'originel est-elle le trait qui définit notre rapport à la nature, ou le symptôme confus d'une exigence de valeur qu'un rapport purement opératoire ne sait pas formuler ?

Trois œuvres permettent d'en juger. Dans \emph{La Connaissance de la vie} (1952, édition revue et augmentée en 1965), Georges Canguilhem, philosophe et médecin, réunit des études d'histoire et de philosophie de la biologie qui se demandent comment la pensée connaît le vivant dont elle procède. \emph{Vingt mille lieues sous les mers} (1869-1870), roman d'aventures scientifiques de la série des « Voyages extraordinaires », fait du \emph{Nautilus} le lieu d'une rencontre entre la machine et l'océan, au siècle qui croit au progrès. \emph{Le Mur invisible} (1963), récit à la première personne écrit par Marlen Haushofer dans l'Autriche de la guerre froide, enferme une femme dans une nature de montagne après une catastrophe inexpliquée.

Nous accorderons d'abord à Dagognet que cette sacralisation de l'originel existe et qu'elle est aveugle. Nous montrerons ensuite qu'elle ne caractérise pourtant pas notre rapport à la nature, qui est d'abord connaissance et transformation, au point que l'opposition même de l'incréé et du fabriqué se défait. Nous soutiendrons enfin que la vraie question n'est pas celle de l'origine, mais celle des normes de vie que nos interventions instituent.

\begin{piege}

L'ouverture cosmique (« De tout temps, l'homme\ldots{} ») ou d'actualité vague (« À l'heure du réchauffement climatique\ldots{} ») : elle ne contient aucun problème et annonce un devoir de culture générale. Le rapport 2026 met aussi en garde contre l'amorce qui consiste « en une autre citation qui parle d'autre chose » (p.~6), et rappelle que « Prendre parti politiquement s'avère par exemple assez maladroit » (p.~1).

\end{piege}

\begin{jury}

Parmi les éléments « incontournables », le rapport 2026 cite « Une présentation des œuvres tenant compte de leur spécificité générique et de leur contexte d'écriture » (p.~7). D'où le paragraphe qui situe les trois œuvres (date, genre, contexte) avant l'annonce du plan. Le jury conseille de partir de références culturelles « qui permettent de montrer d'entrée de jeu que le sujet est bien compris » (p.~6).

\end{jury}

\begin{methode}

L'annonce de plan énonce des thèses, non des thèmes : chaque membre doit être une proposition qu'on pourrait contester.

\end{methode}

\hypertarget{uxe9tape-8-ruxe9diger-le-duxe9veloppement-et-la-conclusion}{%
\subsection{Étape 8 --- Rédiger le développement et la conclusion}\label{uxe9tape-8-ruxe9diger-le-duxe9veloppement-et-la-conclusion}}

\begin{indication}{1}

Chaque paragraphe part d'une idée, la prouve par une ou deux références précises, puis revient explicitement au libellé.

\end{indication}

\begin{indication}{2}

Concédez franchement en première partie : plus la thèse de Dagognet est rendue forte, plus la contestation sera crédible. En troisième partie, appuyez le déplacement sur un concept précis de Canguilhem.

\end{indication}

\begin{indication}{3}

En conclusion, répondez nettement à la question posée, sans citation ni argument nouveaux, et ouvrez à partir d'une limite reconnue du devoir.

\end{indication}

\hypertarget{solution-7}{%
\subsubsection{Solution}\label{solution-7}}

\textbf{I. Une religiosité de l'originel travaille bien notre rapport à la nature, et elle est naïve}

\textbf{A. Le prestige de l'intact.} Que la nature soit spontanément investie d'une aura sacrée, le capitaine Nemo en donne le témoignage le plus éclatant. « La mer est tout », déclare-t-il à Aronnax : elle est l'immense désert où l'homme n'est jamais seul, le souffle pur qui porte la vie, et surtout le seul espace que les despotes ne possèdent pas. Nemo, qui a rompu avec la société tout entière, ne cherche pas dans la mer des ressources mais un salut, et c'est dans l'élément non humain qu'il place la garantie de sa liberté. La virginité des lieux fonctionne comme une valeur en soi : la conquête du pôle Sud, où Nemo plante son pavillon noir, vaut parce que personne n'y est venu. Le début du \emph{Mur invisible} en offre la version profane : la montagne du week-end est un décor choisi pour ce qu'elle a d'intact, et c'est cette intégrité supposée qui la rend désirable. Canguilhem permet de nommer le fond de ce sentiment. Dans « La monstruosité et le monstrueux », il rappelle que le monstre a longtemps été reçu comme un scandale, signe d'une faute ou d'une colère divine : on ne s'indigne d'une infraction que devant une loi tenue pour sacrée. L'horreur du monstrueux atteste donc qu'un ordre naturel est investi d'une autorité quasi religieuse.

\textbf{B. L'oubli des médiations.} Ce qui rend cette dévotion naïve, c'est qu'elle prend pour originel ce qui est déjà travaillé. Nemo prétend fuir les « combinaisons humaines » ; il n'atteint pourtant la mer libre qu'à bord de la machine la plus perfectionnée de son siècle. Le \emph{Nautilus} est un artefact intégral : double coque d'acier, électricité tirée du sodium extrait de l'eau de mer, réservoirs d'air, cuisine et éclairage électriques ; Nemo s'habille d'étoffes de byssus et fume des cigares d'algues ; son salon abrite une bibliothèque de douze mille volumes, des tableaux de maîtres et un orgue. Sa devise, \emph{Mobilis in mobili}, dit qu'on n'habite le mobile qu'en s'en rendant maître. Le Nemo qui célèbre la mer est à l'océan ce que le buveur d'eau « naturelle » est à la source chez Dagognet : un homme qui doit à l'industrie l'accès à ce qu'il croit recevoir de la nature. \emph{Le Mur invisible} aboutit au même constat par la voie inverse. Jetée dans une nature enfin « première », la narratrice ne survit que grâce aux réserves du monde disparu, qu'elle inventorie avec lucidité : allumettes et cartouches comptées, semences de pommes de terre et de haricots, faux, outils, crayons et papier sur lesquels elle rédige son récit. Bella est une vache d'élevage ; le foin se fauche sur des prairies que des générations de paysans ont gagnées sur la forêt. Canguilhem donne à cette naïveté sa forme théorique : dans « Le vivant et son milieu », il retrace l'histoire de la notion de milieu, venue de la mécanique newtonienne, passée à la biologie avec Lamarck, reprise par Comte puis par Darwin. Ce que l'on croit être un donné brut est déjà un concept, avec une histoire ; croire à un pur donné, c'est oublier les médiations par lesquelles on le pense.

\textbf{C. Une nature sans bienveillance.} La religiosité naïve ne croit pas seulement l'originel pur, elle le croit bon. Les œuvres la démentent plus durement encore que Dagognet. La vallée de Haushofer n'a rien d'une mère : le froid, la faim de chaque hiver, l'épuisement, la mort des petits de la chatte, la corneille blanche que sa bande rejette, le gibier qu'il faut tuer pour manger composent une nature indifférente, ce qui est pire, pour qui attendait une consolation, qu'une nature cruelle. Le roman de Verne, sous ses dehors d'émerveillement, dit la même chose : la mer enferme le \emph{Nautilus} sous la banquise jusqu'à l'asphyxie, lance un requin sur le pêcheur de perles de Ceylan, que Nemo affronte au poignard avant que le harpon de Ned Land n'achève la bête, livre un homme d'équipage aux tentacules des poulpes géants, et engloutit le récit dans le Maelström. Canguilhem, enfin, retire à la nature son infaillibilité : la monstruosité atteste que la forme vivante peut manquer son accomplissement, et il ajoute que « la vie est pauvre en monstres », tandis que le fantastique est un monde. Ce que nous redoutons dans la nature est en grande partie notre propre production imaginaire. Cette crainte n'est d'ailleurs pas le contraire de la religiosité, mais son autre face : le sacré fascine et terrifie à la fois, et la mer de Nemo est adorée et redoutée d'un même mouvement. La sacralisation de l'originel repose ainsi sur une double erreur : elle prête au donné une pureté qu'il n'a pas et une bonté qu'il n'a jamais eue.

Dagognet a donc raison de dénoncer une illusion. Mais une illusion suffit-elle à \emph{caractériser} un rapport ?

\textbf{II. Ce n'est pourtant pas le trait qui définit notre rapport à la nature : celui-ci est d'abord opératoire}

\textbf{A. Connaître : la nature comme objet d'analyse.} Avant d'adorer, l'homme classe. \emph{Vingt mille lieues sous les mers} est un roman de savant : Aronnax, professeur au Muséum, auteur des \emph{Mystères des grands fonds sous-marins}, répond à l'inconnu par une hypothèse --- un narval géant --- et non par l'extase ; Conseil récite embranchements, classes et ordres devant chaque poisson qui passe derrière la vitre du salon, tandis que Ned Land n'y voit que du gibier. Le comique de la nomenclature dit une vérité : l'océan devient monde pour Aronnax dans la mesure où il devient catalogue, et le \emph{Nautilus} est autant un musée et un laboratoire qu'un navire. Canguilhem théorise ce geste dans « La pensée et le vivant » : la connaissance cherche la sécurité en réduisant l'inconnu au connu, elle analyse, elle décompose. Mais il en marque aussitôt la limite : la vie est « formation de formes », tandis que la connaissance est analyse des matières informées. « L'expérimentation en biologie animale » en tire les conséquences : l'individualité de l'organisme résiste, le laboratoire est lui-même un milieu qui modifie ce qu'il observe, et l'on ne conclut pas sans précaution d'une espèce à l'autre. Le rapport cognitif est donc premier, mais il ne se referme jamais sur son objet ; ce reste du vivant qui excède l'analyse est peut-être ce que la religiosité naïve traduit maladroitement en langage sacré.

\textbf{B. Transformer : la technique n'est pas le contraire de la nature.} Dagognet a raison : aucun rapport humain à la nature ne se passe de l'outil. Mais il garde l'opposition du donné et du fabriqué, dont il se contente d'inverser les valeurs, louant l'intervention là où ses adversaires louent l'intact. Canguilhem va plus loin dans « Machine et organisme ». Au lieu d'expliquer l'organisme par la machine, à la manière cartésienne, il comprend la machine à partir de l'organisme : la machine est rigide, construite pour une fin extérieure, incapable de se réparer, alors que l'organisme admet suppléances, régulations et usages déviants. Surtout, la technique doit être pensée comme un « phénomène biologique universel », le prolongement de l'activité par laquelle tout vivant organise ce qui l'entoure. Si la technique est un fait de vie, la « combinaison humaine » n'est pas une infraction commise contre la nature, elle en est une production : la ligne de partage sur laquelle repose la religiosité naïve s'efface, mais celle qui permettait à Dagognet d'opposer l'intervention salvatrice à l'âge d'or s'efface avec elle, puisqu'il n'y a jamais eu que des degrés d'intervention. Les deux romans le montrent. Le \emph{Nautilus} vit de l'océan : son électricité vient du sodium marin, son combustible des houillères sous-marines, sa nourriture des bancs de poissons et des algues, sa richesse des trésors engloutis ; il est, au sens de Canguilhem, un organe prolongé greffé sur un milieu. La narratrice de Haushofer redevient paysanne, et le roman consacre de longues pages à la fenaison, à la traite, à la coupe du bois, au potager défendu contre le gibier. Ce travail ne la sépare pas de la nature : c'est lui qui la maintient en vie, physiquement et moralement, car les tâches réglées par les bêtes et les saisons la protègent du désespoir. Là où la religiosité naïve croit qu'il suffit de laisser être, les œuvres montrent qu'il faut faire.

\textbf{C. Habiter : le vivant se fait un milieu.} Reste à comprendre pourquoi cette activité n'est pas une simple domination. « Le vivant et son milieu » distingue, à la suite de von Uexküll, l'environnement géographique (\emph{Umgebung}), le milieu de comportement propre à une espèce (\emph{Umwelt}), fait des seuls signaux qui ont pour elle une valeur, et l'univers de la science (\emph{Welt}). La tique et l'homme, placés dans le même environnement, n'habitent pas le même monde. Le propre du vivant est de centrer son milieu sur lui, d'y introduire des valeurs ; et le milieu du savant, souligne Canguilhem, n'a pas de privilège absolu : il est celui d'un vivant qui a fait de la vérité une de ses valeurs. Les deux fictions le figurent. Le \emph{Nautilus} est un \emph{Umwelt} creusé dans un océan qui, pour l'homme, est d'abord une masse mortelle : derrière la vitre du salon, l'abîme devient spectacle, et cette vitre est l'image même de la culture qui s'interpose entre l'homme et l'élément. \emph{Le Mur invisible} rend la démonstration presque expérimentale : la paroi découpe arbitrairement une vallée, une forêt, un alpage, et l'on voit une femme convertir cet environnement en monde vécu --- chemins tracés, saisons partagées entre le pavillon et l'alpage, calendrier tenu, bêtes nommées. Le mur est un fait ; le domaine est une œuvre. Ni sanctuaire ni matériau, la nature est d'abord, pour le vivant humain, un milieu qu'il habite en le transformant.

Mais si le rapport à la nature est tout entier opératoire, comment distinguer l'intervention qui fait vivre de celle qui détruit ? La conclusion de Dagognet, « Partout, il faut intervenir », ne le dit pas : il manque un critère.

\textbf{III. Le vrai partage n'est pas entre le donné et le fabriqué, mais entre des normes de vie}

\textbf{A. Ce que la préférence pour le donné vise sans savoir le dire.} Une croyance aussi tenace mérite mieux qu'un diagnostic d'ignorance. Dagognet le concède à demi lorsqu'il reconnaît que la solution qu'il rejette « exerce le plus d'attrait » et que Kant « lui-même » a cédé à ce prestige. Or l'exemple kantien dit pourquoi. Ce qui déçoit dans le siffleur caché n'est pas que le chant soit moins beau --- il est identique --- mais qu'il n'engage plus rien : ce n'est plus un vivant qui chante, c'est un garçon qui s'amuse. La préférence pour le naturel n'est donc pas, au fond, une préférence pour l'ancienneté, mais l'exigence que ce que nous admirons existe pour soi et non pour nous. Les romans confirment ce déplacement. Ce qui attache Nemo à la mer n'est pas qu'elle soit vierge, c'est qu'elle soit libre : son culte de l'océan est une politique, celle d'un proscrit qui a tout perdu sous l'oppression. Ce qui lie la narratrice à Lynx, à Bella, à la chatte n'est ni la pureté ni l'antériorité, mais le fait qu'ils sont des existences irremplaçables dont elle a la charge : elle ne les possède pas, elle leur répond. Ce que la religiosité traduit en langage de sacrilège est la conscience d'une fragilité partagée, celle que Canguilhem découvre dans l'angoisse devant le monstre : nous savons obscurément que la vie, en nous aussi, peut échouer.

\textbf{B. Du naturel à la norme.} « Le normal et le pathologique » fournit le critère qui manquait. Le normal n'y est pas la moyenne : la fréquence ne fonde aucune valeur, et l'anomalie n'est pas la maladie. La santé est une marge de tolérance aux infidélités du milieu, la capacité d'instituer de nouvelles normes quand les conditions changent ; la maladie est une vie réduite à une seule norme, rétrécie. La vie est activité normative, et il n'y a « pas de fait normal ou pathologique en soi ». Transposé à notre sujet, ce critère change la question : à « est-ce naturel ? », question d'origine, se substitue « quelles normes de vie cela institue-t-il ? ». Un barrage, une variété sélectionnée, une station d'épuration ne sont ni bons ni mauvais parce qu'ils sont artificiels ; ils le sont selon qu'ils élargissent ou rétrécissent le pouvoir des vivants d'habiter un milieu changeant. Inversement, le retour au sauvage n'est pas bon parce qu'il est premier : la narratrice vit dans une nature intacte, mais sa marge est infime --- une blessure, un hiver plus rude, la dernière allumette, et tout s'effondre. Quant à l'« éco-industrie » de Dagognet, elle ne se justifie pas parce qu'elle est industrielle, mais dans la mesure où elle restitue aux vivants des normes praticables. Verne fournit le contre-exemple qui manque au texte. Nemo intervient partout, avec une puissance sans égale, et rien ne distingue chez lui l'usage du saccage, sinon son bon plaisir : il refuse que Ned Land chasse les baleines, qu'il juge inoffensives, puis fait massacrer un troupeau de cachalots qu'il tient pour cruels. Le partage entre destruction légitime et destruction coupable n'y repose sur aucune norme, mais sur un décret souverain : voilà ce que produit un rapport qui a renoncé à la sacralité du donné sans lui substituer de critère.

\textbf{C. Le désastre ne vient pas de la technique, mais de la violence.} La fin du \emph{Mur invisible} tranche par un fait. Après des années où la forêt n'a jamais fait à la narratrice que ce qu'une forêt fait, un homme surgit sur l'alpage et tue à la hache le jeune taureau et Lynx, sans raison ni profit ; elle l'abat d'un coup de fusil. Ce n'est pas la nature qui détruit le monde qu'elle a construit, c'est le dernier homme. La naïveté n'est pas seulement de croire la nature bonne ; elle est aussi de croire l'homme innocent dès lors qu'il cesse de la sacraliser. Verne conclut de même : le voyage s'achève lorsque Nemo, cessant d'être savant pour redevenir vengeur, éperonne un navire de guerre et le regarde sombrer, avant de s'agenouiller en pleurs devant le portrait d'une femme et de deux enfants ; Aronnax, jusque-là fasciné, prend la fuite. Le \emph{Nautilus} n'est pas condamné par sa technique, mais par la haine qu'il sert. La défiance envers la technique n'est donc pas toujours naïve : la narratrice soupçonne que le mur lui-même est l'effet d'une arme nouvelle, qui a figé les vivants de l'autre côté, et l'éperon du \emph{Nautilus} montre qu'une machine faite pour connaître peut servir à tuer. Mais ce que l'on redoute alors n'est pas l'artifice comme tel : c'est une puissance qui s'affranchit de toute norme de vie. Canguilhem, enfin, rappelle dès son introduction que la pensée du vivant doit tenir du vivant l'idée du vivant : le savoir ne s'exerce pas d'un point de vue extérieur au monde qu'il transforme, et c'est pourquoi il engage une responsabilité.

\textbf{Conclusion}

Il fallait accorder à Dagognet ce qu'il a de solide : une dévotion à l'originel travaille bien notre rapport à la nature, et elle est naïve, puisqu'elle prend pour donné ce qui est déjà cultivé, fauché, sélectionné, et qu'elle prête à la nature une bienveillance que la vallée de Haushofer et la mer de Verne démentent. Mais on ne caractérise pas un rapport par son illusion. Notre rapport premier à la nature est opératoire : il connaît en analysant, il transforme en travaillant, et cette transformation, phénomène vital selon Canguilhem, n'est pas une infraction contre la vie. L'opposition de l'incréé et des « combinaisons humaines » ne résiste donc pas ; seulement, en la retournant au lieu de la quitter, Dagognet laisse son injonction finale sans mesure. Cette mesure, la normativité vitale la fournit : ce qui importe n'est pas l'origine de ce que nous faisons, mais les normes de vie que nous instituons et la marge que nous laissons aux vivants, nous compris. \textbf{La religiosité naïve ne caractérise pas le rapport de l'homme à la nature : elle en est le symptôme, la forme confuse d'une exigence de valeur privée de critère.} Reste à savoir si cette mesure peut se déduire d'un savoir. La narratrice du \emph{Mur invisible}, qui continue d'écrire et de soigner ses bêtes après le meurtre, suggère qu'elle s'apprend plutôt dans la charge que l'on accepte d'un autre vivant.

\begin{piege}

Les citations approximatives entre guillemets. Dans le doute sur la lettre d'un passage, mieux vaut une référence précise sans guillemets (l'épisode, le chapitre, le personnage) qu'une fausse citation.

\end{piege}

\begin{piege}

Le faux dépassement : « il faut trouver un équilibre entre respect de la nature et progrès technique ». Cette phrase ne fournit aucun critère ; le III doit au contraire changer la question.

\end{piege}

\begin{jury}

Selon le rapport 2026, « Les exemples doivent avoir une valeur argumentative et non narrative » (p.~6), et la conclusion doit retracer l'évolution de la réflexion en énonçant « clairement la réponse donnée à la problématique du sujet » (p.~7). Il ajoute que « la longueur de la dissertation n'est pas une fin en soi » (p.~6) : le devoir ci-dessus est dense ; en quatre heures, deux exemples bien exploités par sous-partie suffisent.

\end{jury}

\begin{piege}

Réciter ce corrigé. Le rapport 2026 prévient : « La récitation d'un cours ou d'un corrigé est un symptôme très facilement repérable par le correcteur et très pénalisant pour le candidat. » (p.~5) Retenez les idées et les références, pas les phrases.

\end{piege}

\begin{piege}

Les fautes sur les titres : le jury a lu « Vingt milles lieux » (p.~2). Les pénalités de langue peuvent aller jusqu'à 2 points (p.~2) ; à la main, les titres d'œuvres se soulignent.

\end{piege}

\hypertarget{bilan}{%
\section{Bilan}\label{bilan}}

\begin{itemize}
\tightlist
\item
  \textbf{Résumé} : identifier la polarité du texte (ici, un plaidoyer pour l'intervention contre l'illusion naturaliste) avant toute rédaction ; découper par unités de sens et répartir les mots selon le poids argumentatif ; garder l'ordre et l'énonciation ; reformuler ; fixer une convention de décompte et viser le centre de la fourchette.
\item
  \textbf{Dissertation sur citation} : analyser chaque terme du libellé, en particulier le verbe (\emph{caractériser}) et les termes évaluatifs (\emph{religiosité}, \emph{naïve}, \emph{douteuses}) ; faire apparaître les présupposés de la citation (ici, le partage naturel/artificiel) ; construire un troisième temps qui change la question au lieu de chercher un compromis.
\item
  \textbf{Œuvres} : retenir de Canguilhem quatre outils immédiatement transposables --- la connaissance comme analyse et ses limites, la technique comme phénomène biologique, le milieu propre du vivant, la santé comme normativité --- et de chaque roman une poignée de scènes précises (pour Verne : « La mer est tout », le \emph{Nautilus} nourri de l'océan, Conseil et la nomenclature, les cachalots, l'éperonnage ; pour Haushofer : le pavillon de chasse, l'inventaire des réserves, la fenaison et l'alpage, les bêtes, l'homme à la hache).
\item
  \textbf{Ce qui a départagé les copies en 2026} : le décompte sincère des mots, l'enchaînement des trois raisons dans le résumé, le sens de « religiosité », la place faite à la technique et à Canguilhem, et deux œuvres au moins dans chaque paragraphe.
\item
  \textbf{Questions transposables} : la nature comme norme ou comme matériau, l'opposition du naturel et de l'artificiel, le rôle de la technique dans le rapport au vivant, la responsabilité envers les autres vivants ; le plan « concession --- dissolution de l'opposition --- critère normatif » se réemploie pour la plupart des sujets de la session sur ce thème.
\end{itemize}

\end{document}
